\chapter{Searching across families}
\label{ch:search}

A search draws a new architecture for each trial, trains it for \code{trial_epochs}, and records the score.
The final Adam run happens after the cap is reached, on the single best completed trial, and it starts from a fresh initialization.

\section{Three families, one study}
\label{ex:three-search}

This is the shipped eddy experiment, reduced to a handful of trials so you can see one row from each family.
Copy the ranges from the active \file{user_input.py} and change the budget, the database, and the study name.
The production file uses 500 completed trials, 4000 epochs per trial, and 45000 final Adam epochs.

\begin{lstlisting}
problem_name = "steady_eddies"
constraints = "hard"
model_family = "hybrid_quantum"  # used only when the family is not tuned

model_family_options = (
    "classical",
    "hybrid_quantum",
    "full_quantum",
)
embedding_value = "reupload"
embedding_options = ("reupload",)

use_optuna = True
tune_variables = (
    "model_family",
    "width",
    "depth",
    "n_qubits",
    "q_layers",
)

use_family_specific_search_ranges = True
classical_width_min = 1
classical_width_max = 5
classical_width_step = 1
classical_depth_min = 2
classical_depth_max = 7
hybrid_width_min = 1
hybrid_width_max = 5
hybrid_width_step = 1
hybrid_depth_min = 2
hybrid_depth_max = 7

n_qubits_min = 1
n_qubits_max = 8
q_layers_min = 1
q_layers_max = 8

# Three-tower overlap of the box above. See Table of parameter counts.
total_params_min = 42
total_params_max = 603

lr_value = 1.0e-3
log_w_mass_value = 0.0
log_w_mom_value = 0.0
hybrid_classical_warmup_fraction = 0.0
parameter_objective_mode = "raw_loss"
optimize_total_parameters = True

n_trials = 9
optuna_n_startup_trials = 2000
trial_epochs = 50
trial_nx = 24
trial_nt = 24
score_nx = 32
score_nt = 32

final_adam_epochs = 100
final_adam_nx = 32
final_adam_nt = 32
show_trial_epoch_output = False
show_epoch_output = True
epoch_output_every = 25
use_lbfgs = False

storage_url = "sqlite:///optuna_manual_eddies_three_family.db"
study_name = "manual_eddies_three_family"
\end{lstlisting}

Launch it with one process:

\begin{lstlisting}[style=sh]
python run_pinn.py
\end{lstlisting}

What you should see:

\begin{itemize}
\item A line naming the joint search, the three families, the tune set, \code{n_startup_trials=2000}, and the completed-trial cap.
\item A line saying the search runs until 9 completed trials, each with at least 42 and at most 603 parameters.
\item Discard messages for draws outside that window.
  Those trials are absent from the completed count.
\item One table row per completed trial.
  Classical rows show qubits and layers as 0.
  Full-quantum rows show width and depth as 0.
  Hybrid rows show a positive width, depth, qubit count, and layer count, and the embedding \code{reupload}.
\item Because the startup count is above the cap, families are drawn with equal probability.
  Nine trials will not split exactly three and three and three, but a single family should not take the whole table.
\item After the ninth completion, a \code{Best trial params} listing and a final Adam banner for whichever completed trial had the lowest objective.
\end{itemize}

The learning rate is absent from \code{tune_variables}, so every family uses \(10^{-3}\).
The log-weights are absent, so every family uses PDE weights \((1,1)\).
Width and depth are sampled for classical and hybrid trials and skipped for full-quantum trials.
Qubits and layers are sampled for the two quantum families and skipped for classical trials.

The database is
\file{analysis/optuna_manual_eddies_three_family/optuna_manual_eddies_three_family.db}.
Rerunning the same command continues that study.
If nine trials are already complete, the process adds none and goes straight to final training.
Raise \code{n_trials} to append more, or change both \code{storage_url} and \code{study_name} to start over.

\section{One family, circuit size only}

To tune the circuit of a hybrid model and hold the trunk fixed, omit \code{model_family}, \code{width}, and \code{depth} from the tune list.

\begin{lstlisting}
problem_name = "tidal_test"
constraints = "hard"
model_family = "hybrid_quantum"
model_family_options = ("hybrid_quantum",)
use_optuna = True
tune_variables = ("n_qubits", "q_layers")

width_value = 16
depth_value = 4
lr_value = 1.0e-3
log_w_mass_value = 0.0
log_w_mom_value = 0.0
embedding_value = "reupload"
n_qubits_min = 1
n_qubits_max = 4
q_layers_min = 1
q_layers_max = 3

total_params_min = None
total_params_max = None
n_trials = 6
optuna_n_startup_trials = 6
trial_epochs = 40
trial_nx = 24
trial_nt = 24
score_nx = 32
score_nt = 32
final_adam_epochs = 80
final_adam_nx = 32
final_adam_nt = 32
show_trial_epoch_output = False
use_lbfgs = False

storage_url = "sqlite:///optuna_manual_hybrid_circuit.db"
study_name = "manual_hybrid_circuit"
\end{lstlisting}

Every completed row should say \code{hybrid_quantum}, width 16, and depth 4.
Qubits and layers should vary inside the boxes above.
\code{model_family_options} is unused while the family is not tuned, but a one-element tuple documents the intent and keeps a later edit from accidentally opening the other families inside this study name.
If you later add \code{"model_family"} to \code{tune_variables} on this same \code{study_name}, change the study name first.

The same pattern with \code{model_family = "full_quantum"} and the same tune list searches only the circuit.
Width and depth in the file are ignored.
A classical circuit search is empty: drop \code{n_qubits} and \code{q_layers}, and list \code{"width"} and \code{"depth"} instead.

\section{Classical width and depth}

\begin{lstlisting}
problem_name = "dam_break"
constraints = "soft"
model_family = "classical"
use_optuna = True
tune_variables = ("width", "depth")
use_family_specific_search_ranges = True
classical_width_min = 8
classical_width_max = 24
classical_width_step = 8
classical_depth_min = 2
classical_depth_max = 4

n_trials = 4
optuna_n_startup_trials = 4
trial_epochs = 30
trial_nx = 21
trial_nt = 21
score_nx = 41
score_nt = 41
final_adam_epochs = 60
show_trial_epoch_output = False
use_lbfgs = False
total_params_min = None
total_params_max = None

storage_url = "sqlite:///optuna_manual_classical_search.db"
study_name = "manual_classical_search"
\end{lstlisting}

Widths fall on 8, 16, and 24.
Depths fall on 2, 3, and 4.
The table's qubit and layer columns stay at 0.

\section{Learning rate and loss weights}

Add \code{"lr"} to \code{tune_variables} to sample the Adam rate log-uniformly between \code{lr_min} and \code{lr_max}, for whichever family the trial drew.
Add \code{"log_w_mass"} or \code{"log_w_mom"} to sample the logarithms in Equation~\eqref{eq:weights}.
Sampling either logarithm prints both log-weight columns.
Leaving both at zero and out of the tune list hides those columns, because the weights are then identically \((1,1)\).

Add \code{"embedding"} only when \code{embedding_options} contains more than one layout.
A one-element tuple makes every quantum trial use that layout, and the embedding column stays hidden unless the fixed value is something other than \code{once} or the name is being tuned.
The shipped reupload studies show the column because \code{embedding_value} is \code{reupload}.

\section{Parameter-aware objectives}

To minimize the raw loss, keep the shipped mode:

\begin{lstlisting}
parameter_objective_mode = "raw_loss"
optimize_total_parameters = True
\end{lstlisting}

\code{optimize_total_parameters = False} also returns the raw loss, for any mode.
Counts are still stored.

To charge a model for its size, set

\begin{lstlisting}
optimize_total_parameters = True
parameter_objective_mode = "loss_multiplier"
total_params_objective_weight = 1.0
total_params_reference = 1000.0
\end{lstlisting}

A model with \(N_{\mathrm{ref}}\) parameters is scored as \(2L\).
A model with half that many parameters is scored as \(1.5L\).
Choose the reference near the middle of the window you care about.
With the shipped reference of \(10^{5}\) and a few hundred parameters, the factor is about \(1.006\), so a multiplier of weight 1 barely moves the ranking.
The shipped eddy file avoids that ambiguity by setting the mode to \code{raw_loss}.
In \code{loss_multiplier} mode the weight is live.

To prefer the smallest model that reaches a loss target, first finish a classical baseline and read its \code{raw_objective}.
Then start a new study:

\begin{lstlisting}
optimize_total_parameters = True
parameter_objective_mode = "target_error_min_params"
target_raw_objective = 1.0e-3  # replace with the baseline raw loss
target_error_constraint_weight = 100.0
total_params_reference = 100.0
\end{lstlisting}

Trials above the target are dominated by the excess term.
Trials at or below the target are ranked by \(N/N_{\mathrm{ref}}\).
Use a new \code{study_name}.
The baseline trials were ranked by a different objective, and Optuna's best trial would mix the two scales if they shared a study.

\section{Parallel workers}

\file{run_cores.sh} runs the search as \(N\) processes and then trains the best trial once.

\begin{lstlisting}[style=sh]
./run_cores.sh 4
\end{lstlisting}

\(N\) is a positive integer.
The script creates \file{logs/}, deletes \file{logs/trial_table_rows} so the table header is printed once for this launch, and sets
\code{OMP_NUM_THREADS}, \code{MKL_NUM_THREADS}, \code{TF_NUM_INTRAOP_THREADS}, and \code{TF_NUM_INTEROP_THREADS} to 1.
Quantum trials run a compiled JAX circuit.
One thread per process keeps the workers from oversubscribing the machine.
The script does not pin JAX's own thread pool beyond those variables; a quantum worker can still use more than one core during compilation.

Each worker is \code{python3 run_pinn.py --search-only --worker-id i}, with standard output and standard error in \file{logs/worker_i.txt}.
Completed-trial rows are also written to the terminal.
The workers share the SQLite study.
The first worker creates the database under a file lock.
The storage layer waits up to 120 seconds on a locked database.
Together the workers stop near \code{n_trials} completed trials.
When the parameter window is set, discarded draws do not count, and the workers keep going until the cap of completed trials is reached.

If any worker exits with an error, the script does not start final training.
Read \file{logs/worker_*.txt}.
If every worker succeeds, the script prints \code{Search finished. Training the best completed trial.} and runs \code{python3 run_pinn.py --finalize}.

Do not start a second \code{--finalize}, or a plain \code{python run_pinn.py}, while the workers are still writing the study.
A plain invocation both searches and then trains.
Two final training jobs would write the same weight file.

You can split the same work by hand:

\begin{lstlisting}[style=sh]
python run_pinn.py --search-only --worker-id 0
python run_pinn.py --search-only --worker-id 1
python run_pinn.py --finalize
\end{lstlisting}

Run the two search commands in separate terminals, wait until both have returned, and then finalize.
\code{--search-only} without \code{--worker-id} is an error.

\section{Continuing, extending, and restarting}

The study is reopened by \code{study_name} inside \code{storage_url}.
Completed trials already in the file count toward \code{n_trials}.

\begin{itemize}
\item To add trials, raise \code{n_trials} and run again.
  Leave the family list and the integer boxes unchanged.
\item To change \code{model_family_options}, \code{embedding_options}, or the min, max, or step of a parameter that has already been sampled, create a new \code{study_name} and a new \code{storage_url}.
  Optuna stores the distribution with the study and rejects a conflicting suggestion.
\item To abandon a smoke-test study, point both names at a new database.
  You may delete the old \file{analysis/<name>/} folder when you no longer need its weights.
\item \code{--finalize} never adds trials.
  It trains the best completed trial under the current objective.
  If the study has no completed trial, it raises \code{RuntimeError} and names the study.
\end{itemize}

\section{Which trial is ``best''}

During final training the solver uses Optuna's best trial, the completed trial with the lowest stored objective.
That objective is \(L\) in \code{raw_loss} mode and the transformed value in the other modes.
The analysis script, described in the next chapter, selects the completed trial with the lowest \code{raw_objective} user attribute.
When the mode is \code{raw_loss}, the two selections are the same trial.
When the mode charges for parameter count, Optuna's best trial and the lowest-residual trial can differ.
The summary written by \code{--finalize} records Optuna's best trial.
The analytic report records the lowest residual.
Read the trial number in each file before you compare them.
